% Gerado por scripts/migrate_markdown_to_latex.py.
% Fonte migrada: fases/01_validacao_conceitual/docs/protocolo_experimental.md


\chapter{Protocolo experimental antes dos benchmarks}






Este protocolo define os portoes que uma tecnica precisa atravessar antes de ser usada em comparacoes de desempenho. A intencao e evitar que o trabalho misture tres coisas diferentes: definicao cientifica, validade numerica e ganho real de hardware.



\section{Niveis de certeza}



\begin{longtable}{@{}>{\RaggedRight\arraybackslash}p{0.273\linewidth}>{\RaggedRight\arraybackslash}p{0.273\linewidth}>{\RaggedRight\arraybackslash}p{0.273\linewidth}@{}}
\toprule
Nivel & Pergunta respondida & Evidencia minima \\
\midrule
\endfirsthead
\toprule
Nivel & Pergunta respondida & Evidencia minima \\
\midrule
\endhead
Definicao & O termo esta sendo usado corretamente? & Fonte cientifica ou primaria e definicao operacional. \\
Matematica & A formula local bate com a definicao? & Caso pequeno calculado manualmente e teste deterministico. \\
Algoritmo & A ferramenta implementa o metodo esperado? & Comparacao contra implementacao transparente ou API de referencia. \\
Hardware & A execucao pode sustentar ganho fisico? & Dtype/\allowbreak{}formato/\allowbreak{}kernel/\allowbreak{}medidor mostrando menor memoria, latencia ou maior throughput. \\
\bottomrule
\end{longtable}




Uma tecnica pode passar nos tres primeiros niveis e ainda continuar pendente em hardware. Isso nao e falha; e uma conclusao experimental importante.



\section{Portoes antes de benchmark}


\begin{enumerate}
\item O conceito precisa aparecer na matriz de validacao.
\item A definicao precisa apontar para fonte cientifica ou documentacao primaria.
\item Deve existir teste pequeno e deterministico para o comportamento essencial.
\item A classificacao deve dizer se o resultado e conceitual, numerico,
algoritmico ou hardware.
\item O nucleo da attention deve bater com uma referencia matematica independente
e com APIs publicas de frameworks consolidados.
\item Benchmark amplo so pode ser executado depois dos testes conceituais e da
comparacao entre frameworks passarem.
\end{enumerate}


\section{Validacao cruzada entre frameworks}




Antes dos benchmarks, \texttt{validacao.\allowbreak{}comparacao\_\allowbreak{}frameworks} executa os mesmos tensores \texttt{Q}, \texttt{K} e \texttt{V} em quatro caminhos:


\begin{itemize}
\item formula transparente em NumPy, usada como referencia matematica;
\item formula transparente em PyTorch;
\item \texttt{torch.\allowbreak{}nn.\allowbreak{}functional.\allowbreak{}scaled\_\allowbreak{}dot\_\allowbreak{}product\_\allowbreak{}attention};
\item \texttt{keras.\allowbreak{}ops.\allowbreak{}nn.\allowbreak{}dot\_\allowbreak{}product\_\allowbreak{}attention} com backend TensorFlow.
\end{itemize}






A comparacao usa \texttt{float32}, CPU, dropout desativado, flash attention desativada, seed \texttt{2026}, \texttt{atol=1e-6} e \texttt{rtol=1e-5}. Os casos cobrem valor manual pequeno, multi-head com entradas pseudoaleatorias compartilhadas e mascara aditiva. A entrada e gerada uma unica vez pelo NumPy para evitar divergencia entre os geradores aleatorios dos frameworks.






O CSV desta etapa tem schema proprio e nao altera o schema dos benchmarks de desempenho. A evidencia canônica fica em \texttt{evidencias/\allowbreak{}comparacao\_\allowbreak{}frameworks/\allowbreak{}} e so e aprovada se shape, dtype, finitude e tolerancias passarem em todas as comparacoes.



\section{Recorte experimental}



O nucleo obrigatorio sera um bloco Transformer pequeno e controlado, com:


\begin{itemize}
\item entrada \texttt{[batch, seq\_\allowbreak{}len, d\_\allowbreak{}model]};
\item projecoes densas para Q, K e V;
\item atencao multi-head baseada em \texttt{softmax(QK\textasciicircum{}T /\allowbreak{} sqrt(d\_\allowbreak{}k))V};
\item informacao posicional sinusoidal;
\item residual, normalizacao e feed-forward;
\item saida comparada contra baseline para medir erro.
\end{itemize}




Esse nucleo deve ser chamado de 'bloco Transformer simplificado'. O termo 'Transformer' fica reservado para uma arquitetura completa ou para uma implementacao que tenha justificativa formal para todos os componentes.



\section{Cenarios experimentais}



\begin{longtable}{@{}>{\RaggedRight\arraybackslash}p{0.273\linewidth}>{\RaggedRight\arraybackslash}p{0.273\linewidth}>{\RaggedRight\arraybackslash}p{0.273\linewidth}@{}}
\toprule
Cenario & Descricao & Validade inicial \\
\midrule
\endfirsthead
\toprule
Cenario & Descricao & Validade inicial \\
\midrule
\endhead
\texttt{baseline} & Bloco em \texttt{float32} ou dtype definido como referencia. & Algoritmica \\
\texttt{pruning\_\allowbreak{}magnitude} & Poda por magnitude em pesos lineares. & Conceitual/algoritmica; hardware pendente \\
\texttt{quantization\_\allowbreak{}int8} & Quantizacao int8 validada por dtype, bytes e erro. & Numerica/algoritmica; hardware pendente ate validar kernel \\
\texttt{pruning\_\allowbreak{}plus\_\allowbreak{}quantization} & Combinacao das duas tecnicas, apenas se ambas passarem nos testes isolados. & Pendente \\
\bottomrule
\end{longtable}



\section{Metricas}



As medicoes de qualidade e custo devem ser separadas.



Metricas de qualidade:


\begin{itemize}
\item MSE;
\item MAE;
\item R2;
\item similaridade de cosseno.
\end{itemize}


Metricas de custo:


\begin{itemize}
\item latencia media, p50 e p95;
\item throughput em tokens/s;
\item pico de memoria;
\item FLOPs teoricos.
\end{itemize}



Em GPU, o benchmark deve usar warmup, sincronizacao antes/depois da medicao e registro explicito do dispositivo.



\section{Schema CSV obrigatorio}



Todo resultado experimental deve ser salvo com estas colunas, nesta ordem:



\begin{Verbatim}[fontsize=\small,breaklines=true,label=text]
torch_version,cuda_version,gpu_name,device,model_kind,scenario,batch_size,seq_len,d_model,num_heads,dtype,pruning_sparsity,quantization_method,validity_level,latency_ms_mean,latency_ms_p50,latency_ms_p95,throughput_tokens_s,max_memory_bytes,theoretical_flops,mse,mae,r2,cosine_similarity
\end{Verbatim}




O modulo \texttt{validacao.\allowbreak{}protocolo} valida esse schema antes de qualquer resultado ser aceito como dado experimental do trabalho.



\section{Runner controlado}




O benchmark validado deve ser executado pelo modulo \texttt{validacao.\allowbreak{}benchmark\_\allowbreak{}controlado}. Ele:


\begin{itemize}
\item rejeita cenarios que nao estao registrados no catalogo de validacao;
\item evita APIs depreciadas no pipeline automatizado;
\item fixa seed e registra ambiente em arquivo de metadados;
\item usa o bloco Transformer simplificado como nucleo experimental;
\item mede latencia com warmup, repeticoes e sincronizacao CUDA quando aplicavel;
\item calcula erro contra baseline e FLOPs teoricos;
\item grava CSV apenas se todas as colunas obrigatorias forem validas.
\end{itemize}







Este runner mede o bloco simplificado e deve ser tratado como piloto de infraestrutura. Em CPU, \texttt{max\_\allowbreak{}memory\_\allowbreak{}bytes} e uma estimativa do armazenamento do modelo; apenas em CUDA o campo usa \texttt{torch.\allowbreak{}cuda.\allowbreak{}max\_\allowbreak{}memory\_\allowbreak{}allocated}. A quantizacao manual dequantiza pesos para \texttt{float32}, portanto nao alega armazenamento nem kernel INT8. O benchmark principal das operacoes isoladas e definido separadamente no checklist pre-benchmark.



\section{Runner principal das operacoes isoladas}




\texttt{validacao.\allowbreak{}benchmark\_\allowbreak{}operacoes} implementa o recorte reconciliado entre o plano formal e a reuniao:


\begin{itemize}
\item \texttt{dense\_\allowbreak{}projection}: operacao mantida porque aparece explicitamente no plano;
\item \texttt{self\_\allowbreak{}attention}: uma cabeca com projecoes Q/K/V/saida e SDPA, usada como foco
principal sugerido pelo orientador;
\item o bloco Transformer simplificado fica fora do resultado principal.
\end{itemize}








Os dados sao sinteticos pseudoaleatorios com distribuicao normal, pois o estudo mede operacoes e fidelidade de saida, nao uma tarefa supervisionada ou um dataset. Mesma seed, shapes, entradas e pesos-base sao usados em todos os cenarios. A configuracao schema v2 fixa entradas \texttt{N(0,1)}, pesos com inicializacao Xavier normal (\texttt{desvio=1/\allowbreak{}sqrt(D)} para matrizes quadradas) e bias zero. Essa escala evita que a variancia das projecoes cresca artificialmente com \texttt{D}. Pruning e quantizacao afetam apenas as matrizes de pesos.







A configuracao \texttt{experimentos/\allowbreak{}benchmark\_\allowbreak{}principal\_\allowbreak{}gpu.\allowbreak{}json} codifica seed \texttt{42}, lote \texttt{1}, \texttt{L=\{64,128,256\}}, \texttt{D=\{128,256,512\}}, 20 iteracoes de aquecimento, 50 medicoes e duas execucoes independentes. Cada execucao preserva os mesmos dados e alterna deterministicamente a ordem dos cenarios. O p50 e p95 usam \texttt{numpy.\allowbreak{}percentile} com interpolacao linear.







Cada lote independente grava CSV e JSON de metadados imediatamente. O diretorio tambem recebe snapshot da configuracao, ambiente, log e manifesto com SHA-256. O runner recusa diretorio nao vazio, \texttt{device=auto}, campos nao finitos e benchmark principal sem CUDA/GPU nominal. Assim, uma interrupcao posterior nao apaga o lote ja concluido.







\texttt{validacao.\allowbreak{}analise\_\allowbreak{}benchmark\_\allowbreak{}operacoes} valida os checksums do manifesto antes de ler os dados. Cada candidato e comparado somente com o baseline da mesma configuracao e execucao independente. Depois, as execucoes sao consolidadas com media e desvio-padrao, preservando tambem os resultados individuais. Hashes de entrada e saida precisam coincidir entre execucoes; latencias podem variar.



Comando padrao:



\begin{Verbatim}[fontsize=\small,breaklines=true,label=powershell]
.\.venv\Scripts\python.exe -m validacao.benchmark_controlado --output resultados\benchmark_controlado.csv
\end{Verbatim}



Para uma execucao curta de verificacao:



\begin{Verbatim}[fontsize=\small,breaklines=true,label=powershell]
.\.venv\Scripts\python.exe -m validacao.benchmark_controlado --output resultados\benchmark_teste.csv --warmup 1 --repetitions 3
\end{Verbatim}



\section{Analise dos resultados}





A analise final deve usar \texttt{validacao.\allowbreak{}analise\_\allowbreak{}resultados}. Ela valida o CSV, compara cada cenario com o baseline, gera tabelas de resumo, gera graficos e escreve conclusoes com origem explicita:


\begin{itemize}
\item evidencia no CSV;
\item linha conceitual da matriz;
\item nivel de validade observado.
\end{itemize}



Nenhuma conclusao deve ser escrita como ganho de hardware se o \texttt{validity\_\allowbreak{}level} do resultado nao for \texttt{hardware}.



\section{Validacao da NN como Transformer}





A rede neural usada no nucleo experimental deve ser auditada por \texttt{validacao.\allowbreak{}auditoria\_\allowbreak{}transformer} antes dos benchmarks. O resultado aceito para o recorte atual e \texttt{bloco\_\allowbreak{}transformer\_\allowbreak{}simplificado}.





Uma NN comum com camadas lineares nao pode passar nessa auditoria. Isso e um controle negativo obrigatorio para garantir que o estudo nao esta chamando uma rede neural generica de Transformer.



\section{Comparacao com literatura}



Nesta fase, a comparacao com literatura deve ser feita assim:


\begin{itemize}
\item formulas e definicoes: comparacao direta com valores derivados;
\item bibliotecas: comparacao contra implementacao transparente ou documentacao
primaria;
\item desempenho: comparacao apenas como tendencia ou ordem de grandeza, salvo se
modelo, hardware, kernel, dtype, entrada e protocolo forem equivalentes.
\end{itemize}


\section{Criterio para conclusao de hardware}




Uma tecnica so pode receber nivel 'hardware' quando houver evidencia de que o caminho executado usa a representacao otimizada:


\begin{itemize}
\item quantizacao: armazenamento menor e kernel/caminho compativel com o dtype;
\item pruning: formato ou kernel que explore esparsidade;
\item KV cache: reutilizacao de K/V entre passos e equivalencia numerica com a
execucao sem cache;
\item memoria: medicao real ou estimativa teorica claramente separada;
\item latencia/throughput: medicao com repeticoes, warmup e sincronizacao.
\end{itemize}
